
Vision Foundation Models (VFMs) such as DINOv2~\cite{dinov2} and CLIP~\cite{radford2021clip} have emerged as powerful tools for extracting semantic features. However, their utility in dense prediction tasks is often limited by the low resolution of their output feature maps. Recent advancements in ``universal feature upsampling,'' specifically the state-of-the-art model AnyUp, have addressed this by enabling continuous, backbone-agnostic upsampling. Despite its success with static images, AnyUp operates on a frame-by-frame basis when applied to video, ignoring temporal correlations. This independence results in severe temporal inconsistency, manifesting as flickering artifacts and unstable object boundaries in video sequences.

This thesis proposes ``AnyUp3D,'' a novel extension of the AnyUp framework designed to introduce temporal consistency while preserving the model's universality. The primary objective is to bridge the gap between image-based universal upsamplers and task-specific video processing methods. The proposed methodology extends the upsampling pipeline with 3D convolutions in the encoder to capture spatiotemporal patterns, and 3D cross-attention with temporal masking to correct frame-level inconsistencies by attending to neighboring frames within a spatiotemporal window.

The system is designed to be trained on the UCF-101 dataset to learn robust temporal priors without being tied to a specific feature extractor. To validate the model's effectiveness, we evaluate its performance on the task of Video Object Segmentation using the DAVIS-2017 dataset. AnyUp3D demonstrated significant improvements in temporal stability, measured by cosine similarity between adjacent frames, compared to frame-by-frame baselines, while maintaining high spatial fidelity. This work contributes to the field by providing the first task-agnostic, temporally consistent feature upsampler\footnote{The closest prior methods differ in key respects: FeatUp~\cite{featup} and LIIF~\cite{liif} operate on single frames and have no temporal modeling; video super-resolution methods (e.g., Upscale-A-Video~\cite{upscale_video}) are task-specific and operate in the pixel domain rather than the feature domain. AnyUp3D is the first upsampler that is simultaneously backbone-agnostic, resolution-agnostic, and temporally consistent.}, enabling the seamless application of powerful image-based foundation models to complex video understanding tasks.